\section{A boundary bilinear identity and the exact finite trace coefficient}
\label{sec:spectral-boundary-flux}

\begin{lemma}[AC-S43-flux: an exact translated boundary pairing]
\label{AC-S43-flux}
Fix \(I=(0,\ell)\) and the actual Weil-window operator \(A\).
For real \(u,w\in\mathcal X_I\), define
\[
 F_h(u,w)=\int_h^\ell
        \bigl((Au)(x)w(x-h)-u(x)(Aw)(x-h)\bigr)dx,
                  \qquad 0<h<\ell/4.
                                                        \tag{BT.1}
\]
Then
\[
 \lim_{h\downarrow0}\frac{F_h(u,w)}h
       =\tau_+(u)\tau_+(w)-\tau_-(u)\tau_-(w).           \tag{BT.2}
\]
In particular, if \(Au=\alpha u\) and
\((A-\alpha)w=g\in C_c^\infty(I;\mathbb R)\), then
\[
 \int_Iu(x)g'(x)dx
       =\tau_+(u)\tau_+(w)-\tau_-(u)\tau_-(w).           \tag{BT.3}
\]
This identity does not assume that either eigenfunction has an
\(L^2\) derivative or that the ordinary endpoint derivative exists.
\end{lemma}
\begin{proof}
\proofstep{1. A full-line operator with the exact part.}
By (RN.4) and (EP.6), choose a finite real symmetric signed
measure \(\sigma\) on \(\mathbb R\) such that
\[
 H=G+\sigma*,\qquad
 G=\tfrac12\log(1+D^2),\qquad A=PH E.                   \tag{BT.4}
\]
Here \(E\) denotes zero extension and \(P\) restriction to \(I\).
For example \(\sigma\) contains the scalar atom,
the integrable kernel \(-k\), the finitely many prime atoms,
and the polar kernel cut off at \(|x|\leq\ell\).
That cutoff does not change its compression to \(I\).
The operator \(H\) is self-adjoint on \(\operatorname{Dom}(G)\)
and commutes with whole-line translations \(S_h\).

The Carleman domain splitting proved in AC-S37 and AC-S40 gives
\(Eu,Ew\in\operatorname{Dom}(G)\).
Consequently the following equality is a legitimate \(L^2\)
operator pairing:
\[
                   \langle HEu,S_hEw\rangle
                        =\langle Eu,S_hHEw\rangle.        \tag{BT.5}
\]
The singular part of \(HEu\) outside \(I\) is
\(-\int_I n(x-y)u(y)dy\), where
\(n(t)=e^{-|t|}/(2|t|)\).
The remaining convolution is bounded there by
\(\|\sigma\|_{\rm TV}\|u\|_\infty\).
The same assertions hold for \(w\).

\proofstep{2. The exterior convolution has a precise leading term.}
For \(s>0\), write
\[
 C_u^+(s)=\int_0^\ell n(s+t)u(\ell-t)dt,
 \qquad
 C_u^-(s)=\int_0^\ell n(s+t)u(t)dt.
                                                        \tag{BT.6}
\]
The endpoint trace of AC-S37 implies
\[
 \frac{C_u^\pm(s)}{\sqrt{\log(1/s)}}\longrightarrow\tau_\pm(u),
 \qquad |C_u^\pm(s)|\leq C\sqrt{\log(1/s)}
                              \quad(s\hbox{ sufficiently small}).
                                                        \tag{BT.7}
\]
We check the constant.  Near the relevant endpoint,
\(u(t)=\tau(\log(1/t))^{-1/2}
       +o((\log(1/t))^{-1/2})\).
Fix \(\delta<\min(\ell/2,e^{-2})\).
Replacing \(e^{-(s+t)}/(2(s+t))\) by \(1/(2(s+t))\)
changes its integral by \(O(1)\).
The part \(0<t<s\) is \(O((\log(1/s))^{-1/2})\), by
the boundary upper bound and monotonicity of its logarithmic profile.
For \(s<t<\delta\), replacement of \(1/(s+t)\) by \(1/t\)
has bounded error, since
\(s\int_s^\delta |u(t)|t^{-2}dt\leq\|u\|_\infty\).
The leading contribution is therefore
\[
 \frac{\tau}{2}\int_s^\delta
                     \frac{dt}{t\sqrt{\log(1/t)}}
                           =\tau\sqrt{\log(1/s)}+O(1).
                                                        \tag{BT.8}
\]
An \(o(1/\sqrt{\log(1/t)})\) remainder contributes
\(o(\sqrt{\log(1/s)})\): bound it by an arbitrarily small
multiple of the displayed profile near zero and keep the rest
of the integral fixed.  The part beyond \(\delta\) is bounded.
This proves both assertions in (BT.7), including their coefficient.

\proofstep{3. The two shrinking strips.}
Split (BT.5) into the overlap \((h,\ell)\) and the two remaining
strips.  With the signs in the exterior action just recorded,
one obtains
\[
 \begin{split}
 F_h(u,w)={}&\int_0^h C_u^+(s)w(\ell-h+s)ds
              -\int_0^h u(x)C_w^-(h-x)dx+E_h,\\
                    &E_h=o(h).
 \end{split}                                             \tag{BT.9}
\]
For the error assertion, every bounded convolution contribution
is bounded by a fixed constant times the integral of \(|u|\)
or \(|w|\) on an endpoint strip.
The continuous zero boundary values, or the upper bound (FE.8),
make this \(o(h)\).

Put \(s=ht\) in the first integral.  By (BT.7) and (RN.1),
its integrand tends pointwise on \(0<t<1\) to
\(\tau_+(u)\tau_+(w)\).
For small \(h\), its absolute value is bounded by
\[
 C\left(\frac{\log(1/h)+|\log t|}
                   {\log(1/h)+|\log(1-t)|}\right)^{1/2}
                          \leq C\sqrt{1+|\log t|}.       \tag{BT.10}
\]
The last function is integrable on \((0,1)\).
For the second integral, the same substitution gives the
integrable majorant \(C\sqrt{1+|\log(1-t)|}\) and pointwise
limit \(\tau_-(u)\tau_-(w)\).
Dominated convergence in (BT.9) proves (BT.2).

Finally suppose the hypotheses for (BT.3) hold.
Self-adjointness gives \(\int_Iug=0\).
The two eigenvalue terms cancel exactly in (BT.1), so
\(F_h(u,w)=-\int_h^\ell u(x)g(x-h)dx\).
For sufficiently small \(h\), compact support of \(g\) permits
integration over all of \(I\).
Taylor's formula for the fixed smooth \(g\) gives
\[
 F_h(u,w)=h\int_Iu(x)g'(x)dx+O(h^2)\|u\|_1.
\]
Combine this with (BT.2) to obtain (BT.3).
\end{proof}

\begin{theorem}[AC-S43: the finite coefficient equals the renewal trace]
\label{AC-S43}
For every fixed window with a simple even lowest eigenline,
the coefficient of AC-S42 is exactly
\[
                      \kappa_\lambda=\tau_+(u)=\tau_-(u)
                                      \ne0.              \tag{BT.11}
\]
Consequently
\[
 \sqrt{\log N}\,u_N(\ell)\longrightarrow\tau_+(u),
 \qquad
 \frac{\det_{\rm reg}(D^{(\lambda,N)}_{\log}-z)}
                           {\sqrt{\log N}}
       \longrightarrow-i\lambda^{-iz}
                              \frac{\widehat u(z)}{\tau_+(u)}
                                                        \tag{BT.12}
\]
locally uniformly for the determinant, with the corresponding
\(L^2\) and form-norm vector convergence.
The trace nonvanishing condition of AC-S39 is automatic under
the simple-even minimum hypothesis.
On the unconditional positive range of AC-S38 this trace is
positive, so \(u_N(\ell)>0\) for all sufficiently large \(N\)
in the chosen phase.
\end{theorem}
\begin{proof}
Choose the real odd compact smooth forcing \(g\) of AC-S41,
with \(w=R_og\) and \(\tau_+(w)\ne0\).
Reflection gives
\(\tau_-(u)=\tau_+(u)\) and
\(\tau_-(w)=-\tau_+(w)\).
Equation (BT.3) becomes
\[
                    \int_Iu g'
                         =2\tau_+(u)\tau_+(w).
\]
Comparison with the exact coefficient formula (FA.2) gives
\(\kappa_\lambda=\tau_+(u)\).
The nonvanishing is already proved in AC-S42, so no boundary
unique-continuation assertion is being inserted here.
The remaining conclusions are (FW.9), (FA.4), and AC-S38.
For completeness, the smoothed normalizers \(b_N\) of (RN.22)
and the unsmoothed ones now obey
\(b_N/(\sqrt{\log N}\,u_N(\ell))\to1\), since both numerator
and denominator converge to this same nonzero trace.
\end{proof}

\begin{corollary}[AC-S43-trace-rank: endpoint traces on an eigenspace]
\label{AC-S43-trace-rank}
On every real eigenspace of \(A\),
\[
                \tau_+(u)\tau_+(w)=\tau_-(u)\tau_-(w)
                             \quad\hbox{for all }u,w.    \tag{BT.13}
\]
Thus either both endpoint functionals vanish on that eigenspace,
or they agree up to one fixed sign.  The combined endpoint trace
has rank at most one.  This does not prove simplicity: injectivity
of that trace on a general eigenspace has not been established.
\end{corollary}
\begin{proof}
If \(Au=\alpha u\) and \(Aw=\alpha w\), the integrand in
(BT.1) is identically zero.  This proves (BT.13).
If one functional is nonzero, choose \(u\) on which it is nonzero.
Equation (BT.13) with \(w=u\) makes the two values equal up to
a sign; with this fixed \(u\) and arbitrary \(w\) it gives
the same sign for the two functionals.  Otherwise both are zero.
The rank assertion follows, without a claim about their kernel.
\end{proof}
